\input{../tex-inputs/latex-leseansicht-vorspann}

\section[Gustav Schwarzkopf an Arthur Schnitzler, 19. 1. 1910]{L04267 Gustav Schwarzkopf an Arthur Schnitzler, 19. 1. 1910}
\nopagebreak\mylabel{L04267v}
\rehead{ }\normalsize\beginnumbering\briefempfaengerindex{Schnitzler, Arthur@\textsc{Schnitzler, Arthur}!zzzSchwarzkopf, Gustav@\emph{von Gustav Schwarzkopf}!1910-01-191@{19. 1. 1910}|(be}
\toendnotes[C]{\smallbreak\pagebreak[2]}
\correspDesc{Versand  durch Gustav Schwarzkopf am 19. 1. 1910 in Wien
\newline{}Zustellung  im Zeitraum 19. 1. 1910 in Wien
\newline{}Erhalt  durch Arthur Schnitzler am 24. 1. 1910 in Wien}\toendnotes[C]{\smallbreak}
\Standort{CUL, Schnitzler, B 96a.}
\physDesc{Briefkarte, 531 Zeichen
\newline{}Handschrift: schwarze Tinte, deutsche Kurrent}\toendnotes[C]{\smallbreak}
\pstart
           \raggedleft{}{\pb}19. I 10\pend
           \vspace{0.5em}
\pstart
           Lieber Arthur, ich übermittle Ihnen
      den beifolgenden Brief, den ich eben erhalten.
      Ich tue das nur, weil Sie ja doch nicht
      verſchont bleiben würden; we{\geminationn} ich die
               Vermittlung ablehne, würde H.\pwindex{Holzer, Rudolf 28.\,7.\,1875 Wien – 17.\,7.\,1965 ebd.@\textsc{Holzer, Rudolf} (28.\,7.\,1875 Wien – 17.\,7.\,1965 ebd.), \emph{Schriftsteller, Journalist}|pw} –
      wie aus{ }ſeinem Briefe erſichtlich – Ihnen{ }ſelbſt{ }ſchreiben. We{\geminationn} Sie in der
      Sache etwas tun wollen – bitte,
      \uline{nur}\uline{da{\geminationn}}, we{\geminationn}{ }ſie \uline{Ihnen} einer
      Berückſichtigung wert{ }ſcheint,
      nicht etwa aus Gefälligkeit
            {\pb}für mich, de{\geminationn} ich{ }ſtelle keine
      Bitte und möchte Ihrem
         Bruder\pwindex{Schnitzler, Julius 13.\,7.\,1865 Wien – 29.\,6.\,1939 ebd.@\textsc{Schnitzler, Julius} (13.\,7.\,1865 Wien – 29.\,6.\,1939 ebd.), \emph{Chirurg}|pwv} gegenüber nicht einmal
      als Vermitter gelten.\pend
           
\pstart
           Ihre Frau\pwindex{Schnitzler, Olga 17.\,1.\,1882 Wien – 13.\,1.\,1970 Lugano@\textsc{Schnitzler, Olga} (17.\,1.\,1882 Wien – 13.\,1.\,1970 Lugano), \emph{Schauspielerin, Sängerin}|pwv} und Sie herzlichſt{\\[\baselineskip]}grüßend{\\[\baselineskip]}Ihr{\\[\baselineskip]}\spacefill\mbox{Gustav S}\pend
           \leftskip=0em{}\selectlanguage{ngerman}\endnumbering\briefempfaengerindex{Schnitzler, Arthur@\textsc{Schnitzler, Arthur}!zzzSchwarzkopf, Gustav@\emph{von Gustav Schwarzkopf}!1910-01-191@{19. 1. 1910}|)be}\mylabel{L04267h}
\begin{anhang}
\end{anhang}\newcommand{\dateiname}{L04267}\newcommand{\titel}{Gustav Schwarzkopf an Arthur Schnitzler, 19. 1. 1910}\newcommand{\editorInnen}{Herausgegeben von Jahnke, SelmaMüller, Martin Anton}\input{../tex-inputs/latex-leseansicht-abspann}
